% Generated draft: requires mathematical and source review.
\documentclass[11pt]{article}
\usepackage[margin=25mm]{geometry}
\usepackage{fontspec}
\setmainfont{Latin Modern Roman}
\usepackage{amsmath,amssymb,mathtools,unicode-math}
\setmathfont{Latin Modern Math}
\usepackage{enumitem,needspace,xurl,hyperref}
\setlength{\parindent}{0pt}
\setlength{\parskip}{5pt}
\setlength{\emergencystretch}{3em}
\newcommand{\N}{\mathbb{N}}
\newcommand{\Z}{\mathbb{Z}}
\newcommand{\Q}{\mathbb{Q}}
\newcommand{\R}{\mathbb{R}}
\newcommand{\C}{\mathbb{C}}
\newcommand{\F}{\mathbb{F}}
\newcommand{\A}{\mathbb{A}}
\newcommand{\PP}{\mathbb{P}}
\newcommand{\E}{\mathbb{E}}
\newcommand{\eps}{\varepsilon}
\newcommand{\dd}{\,\mathrm{d}}
\begin{document}
\section*{UnsolvedMath ID: 800539}
\textbf{Source problem code:} catalogue entry
\subsection{Definitions and mathematical statement}
Let \(G=(V(G),E(G))\) be a finite simple undirected graph, and let \(|V(G)|\) denote its number of vertices. This graph convention is the intended interpretation of the extremal question; it is not explicitly specified in the supplied excerpt. Write \(\operatorname{obs}_{\mathrm{src}}(G)\) for the obstacle number of \(G\) under the conventions of the cited source. The supplied record does not define obstacle number or specify admissible obstacles, representations, visibility boundary conventions, or general-position requirements. Thus the formulation below is conditional on recovering that definition from the source; those geometric conventions are not fixed by this draft.

For each integer \(h\geq 2\), define
\[
\mathcal{G}_h=\{G: G\text{ is a finite simple undirected graph and }\operatorname{obs}_{\mathrm{src}}(G)=h\}.
\]
Set
\[
m(h)=
\begin{cases}
\min\{|V(G)|:G\in\mathcal{G}_h\},&\mathcal{G}_h\neq\varnothing,\\
+\infty,&\mathcal{G}_h=\varnothing.
\end{cases}
\]
The value \(+\infty\) is only a convention making the extremal quantity defined if the relevant class is empty.

\textbf{Problem.} Determine \(m(h)\) for every integer \(h\geq 2\). The requirement is obstacle number exactly \(h\).

The supplied survey excerpt reports that Alpet et al. [akl-ong-2010] and Mukkamala et al. [mps-glon-2010] studied obstacle numbers \(1\) and \(2\), and attributes to Mukkamala et al. the existence, for every integer \(h\geq 3\), of a graph with obstacle number exactly \(h\). Under the finite-graph interpretation, this reported existence statement makes \(m(h)\) finite for \(h\geq 3\). The excerpt does not explicitly supply the corresponding existence statement for \(h=2\). It attributes the extremal question to Alpet et al. and describes it as still open. These are supplied source assessments, not independently verified claims.
\subsection{Short English statement}
For each integer \(h>1\), determine the fewest vertices needed by a graph whose obstacle number is exactly \(h\). The source's definition of obstacle number must be recovered before the geometric formulation is complete. The supplied record labels the problem open; this draft does not independently verify that status.
\subsection{Sources}
\begin{itemize}
\item Local UnsolvedMath registry, permanent ID \(800539\), source problem code \texttt{catalogue entry}, title ``Visibility Graphs of Points, Segments and Polygons --- Problem 9.'' The supplied record lists Ghosh and Goswami as proposers, gives the year 2010, and carries the status \texttt{open}. Its creation and update timestamps are both \texttt{2026-10-16T14:00:00Z}; these are registry metadata, not independent mathematical verification.
\item Ghosh and Goswami, \emph{Unsolved Problems in Visibility Graphs of Points, Segments and Polygons} (2010), Problem 9. Supplied source URL: \url{https://arxiv.org/abs/1012.5187}. The provenance records access on 2026-10-16 and extraction by \texttt{arxiv\_latex}. This access date is supplied provenance; no source visit or independent verification is asserted here.
\item Alpet et al., source reference label [akl-ong-2010]. The supplied excerpt attributes the extremal question and work on obstacle numbers \(1\) and \(2\) to this reference. The spelling ``Alpet'' is retained from the record. No full bibliographic entry or separate URL is supplied.
\item Mukkamala et al., source reference label [mps-glon-2010]. The supplied excerpt attributes work on obstacle numbers \(1\) and \(2\), and the existence result for every integer \(h\geq 3\), to this reference. No full bibliographic entry or separate URL is supplied.
\end{itemize}
\end{document}
